\subsection{Пружинки с трением и без}
%%%%%%%%%%%%%%%%%%%%%%%%%%%%%%%%%%%%%%%%%%%%%%%%%%%%
%%%%%%%%%%%%%%%%%%%%%%%%%%%%%%%%%%%%%%%%%%%%%%%%%%%%	
\z[{\begin{wrapfigure}[5]{r}{0.3\textwidth}%
		\vspace{-10pt}%
		\begin{flushright}%
 			\includegraphics[width=0.3\textwidth]{7.7.1}%
		\end{flushright}%
	\end{wrapfigure}%
	}]%
	{На наклонённой под углом $\alpha$ ($\cos\alpha = \frac{5}{7}$) к горизонту поверхности лежит брусок, прикреплённый к упругой невесомой и достаточно длинной пружине. Коэффициент трения бруска о поверхность $\mu = \frac{7}{30}$. Брусок отклоняют вниз вдоль поверхности на расстояние $A_0 = 32$ см от точки $O$, соответствующей положению равновесия бруска при отсутствии трения. Брусок отпускают, и начинаются затухающие колебания. Если брусок подвесить на этой пружине, то она удлиняется на $x_0 = 30$ см.
	\lpu На каком расстоянии от точки $O$ окажется брусок при первой остановке?\lpuend
	\lpu На каком расстоянии от точки $O$ брусок остановится окончательно? \lpuend
	}
%
{\puan  $A_1 = A_0 - 2\mu x_0 \cos\alpha = 27\,\text{см};$~\puan на $3\,\text{см}$ ниже точки $O$}
%%%%%%%%%%%%%%%%%%%%%%%%%%%%%%%%%%%%%%%%%%%%%%%%%%%% 
\z[{\begin{wrapfigure}[6]{r}{0.3\textwidth}%
		\vspace{-10pt}%
		\begin{flushright}%
 			\includegraphics[width=0.3\textwidth]{7.7.2}%
		\end{flushright}%
	\end{wrapfigure}%
	}]%
	{С наклонной плоскости, образующей угол $\alpha$ с горизонтом, начинают соскальзывать два тела массы $m$ каждое, связанные первоначально недеформированной пружиной жесткости $k$. Определите наибольшее удлинение пружины, если трением между нижним телом и плоскостью можно пренебречь, а коэффициент трения между верхним телом и плоскостью равен $\mu$.}
%
{если $\mu < \tg\alpha$, то $x = \frac{mg\mu \cos\alpha}{k}$;
\par
если $\tg\alpha < \mu < \tg3\alpha$, то $x = \frac{mg\mu \cos\alpha}{2k}\left(1 + \sqrt{1 - 2\left(1 - \frac{\tg\alpha}{\mu}\right)^2}\right)$;
\par
если $\tg3\alpha < \mu$, то $x = \frac{2mg \sin\alpha}{k}$
}
%%%%%%%%%%%%%%%%%%%%%%%%%%%%%%%%%%%%%%%%%%%%%%%%%%%%
\z[{\begin{wrapfigure}[5]{r}{0.3\textwidth}%
		\vspace{-10pt}%
		\begin{flushright}%
 			\includegraphics[width=0.3\textwidth]{7.7.3}%
		\end{flushright}%
	\end{wrapfigure}%
	}]%
	{Тележка находится на горизонтальной поверхности стола. На шероховатой горизонтальной поверхности тележки находится брусок, прикреплённый к тележке лёгкой упругой пружиной. Масса тележки в два раза больше массы бруска. Брусок отклоняют влево так, что пружина сжата на величину $x$, а тележка прижата к упору. Затем брусок отпускают. Найдите: 
	\lpu деформацию пружины в момент отрыва тележки от упора; \lpuend
	\lpu скорость бруска в момент отрыва тележки от упора;\lpuend
	\lpu скорость тележки после прекращения движения по ней бруска. \lpuend
	
	\noindent
	Известно следующее. Если брусок подвесить на пружине, то деформация пружины равна $3x$. Если брусок тащить по неподвижной тележке с постоянной скоростью, прикладывая горизонтальную силу к прикреплённой к бруску пружине, то деформация пружины равна $\frac{3x}{4}$. 
	\newpage
	}
%
{\puan  $x_1 = \frac{3x}{4};$~\puan  $v = \sqrt{\frac{3gx}{12}};$~\puan $v = \sqrt{\frac{3gx}{36}}$}
%%%%%%%%%%%%%%%%%%%%%%%%%%%%%%%%%%%%%%%%%%%%%%%%%%%%	
\z[{\begin{wrapfigure}[6]{r}{0.3\textwidth}%
		\vspace{-0pt}%
		\begin{flushright}%
 			\includegraphics[width=0.3\textwidth]{7.7.4}%
		\end{flushright}%
	\end{wrapfigure}%
	}]%
	{Тележка соединена со стеной жёстким стержнем. К её упору прикреплена пружина, другой конец которой связан с бруском (см. рисунок). Вначале пружина не деформирована. На брусок в течение некоторого времени действует постоянная горизонтальная сила $F$, направленная вдоль тележки. После прекращения действия этой силы брусок ещё некоторое время смещается в сторону упора и возвращается, остановившись в исходной точке. Сила трения, действующая со стороны тележки на брусок, равна $f$. Найдите: 
	\lpu с какой силой $N$ тележка давила на стержень в момент прекращения действия силы $F$;\lpuend
	\lpu наибольшее значение силы $N_{\max}$ давления тележки на стержень. \lpuend
	}
%
{\puan  $N = f\left(1 + \frac{4f}{F}\right);$~\puan $N_{\max} = 3f$}
%%%%%%%%%%%%%%%%%%%%%%%%%%%%%%%%%%%%%%%%%%%%%%%%%%%% 
\z[{\begin{wrapfigure}[5]{r}{0.3\textwidth}%
		\vspace{-5pt}%
		\begin{flushright}%
 			\includegraphics[width=0.3\textwidth]{7.7.5}%
		\end{flushright}%
	\end{wrapfigure}%
	}]%
	{На гладкой горизонтальной поверхности тележки находится брусок, прикреплённый к тележке лёгкой упругой пружиной жёсткостью $k = 10$ Н/м (см. рисунок). Тележка с бруском движутся со скоростью $u = 0.25$ м/с по горизонтальной поверхности пола. Пуля, летящая горизонтально со скоростью $33u$, попадает в брусок и застревает в нём. Массы пули и бруска $m = 0.5$ г и $7m$, масса тележки намного больше массы бруска. Направления всех движений находятся в одной вертикальной плоскости. Найдите: 
	\lpu скорость бруска $v_1$ относительно тележки сразу после попадания пули;\lpuend
	\lpu скорость бруска $v_2$ относительно пола сразу после попадания пули; \lpuend
	\lpu максимальную деформацию пружины.\lpuend
	}
%
{\puan  $v_1 = 4u = 1\,\text{м/с};$~\puan  $v_2 = 5u = 1.25\,\text{м/с};$~\puan $x = 8u\sqrt{\frac{2m}{k}} = 2\,\text{см}$}